\documentclass[11pt,letterpaper]{article}
% Reusable style. Build each .tex from its own directory using scripts/build.sh.
\usepackage[T1]{fontenc}
\usepackage[utf8]{inputenc}
\usepackage{lmodern}
\usepackage[margin=0.86in,headheight=15pt]{geometry}
\usepackage{amsmath,amssymb,amsthm,mathtools}
\usepackage{booktabs,longtable,tabularx,array,multicol}
\usepackage[table]{xcolor}
\usepackage{microtype}
\usepackage{enumitem}
\usepackage{fancyhdr}
\usepackage{titlesec}
\usepackage{hyperref}
\usepackage{listings}
\usepackage[most]{tcolorbox}
\usepackage{needspace}
\definecolor{aeNavy}{HTML}{172D43}
\definecolor{aeBlue}{HTML}{176B90}
\definecolor{aeTeal}{HTML}{148477}
\definecolor{aePale}{HTML}{F2F7F9}
\definecolor{aeGray}{HTML}{5C6771}
\hypersetup{colorlinks=true,linkcolor=aeBlue,urlcolor=aeBlue,pdfauthor={Algorithm Engineering Documentation Suite},pdfsubject={Independent educational materials based on the supplied TAOCP volume topics}}
\titleformat{\section}{\Large\bfseries\color{aeNavy}}{\thesection}{0.6em}{}
\titleformat{\subsection}{\large\bfseries\color{aeBlue}}{\thesubsection}{0.6em}{}
\titleformat{\subsubsection}{\normalsize\bfseries\color{aeNavy}}{\thesubsubsection}{0.6em}{}
\setlist[itemize]{leftmargin=*,topsep=3pt,itemsep=1pt}
\setlist[enumerate]{leftmargin=*,topsep=3pt,itemsep=3pt}
\setlength{\parindent}{0pt}\setlength{\parskip}{5pt}
\pagestyle{fancy}\fancyhf{}\fancyhead[L]{\footnotesize\color{aeGray}ALGORITHM ENGINEERING}\fancyhead[R]{\footnotesize\color{aeGray}DOCUMENTATION SUITE}\fancyfoot[L]{\footnotesize\color{aeGray}Independent study and engineering practice}\fancyfoot[R]{\footnotesize\thepage}
\renewcommand{\headrulewidth}{0.35pt}
\newtcolorbox{engineeringnote}[1][]{breakable,colback=aePale,colframe=aeBlue,boxrule=0.5pt,arc=1mm,fonttitle=\bfseries,title={Engineering note},#1}
\newtcolorbox{evidencebox}[1][]{breakable,colback=white,colframe=aeTeal,boxrule=0.6pt,arc=1mm,fonttitle=\bfseries,title={Evidence and deliverables},#1}
\lstdefinestyle{algo}{basicstyle=\ttfamily\footnotesize,breaklines=true,frame=single,rulecolor=\color{aeGray!40},backgroundcolor=\color{aePale},showstringspaces=false,tabsize=2,columns=fullflexible,keepspaces=true,keywordstyle=\color{aeBlue}\bfseries,commentstyle=\color{aeTeal}}
\lstset{style=algo}
\newcommand{\DocTitle}[3]{%
\hypersetup{pageanchor=false}\begin{titlepage}\vspace*{1.2cm}
{\color{aeTeal}\rule{\linewidth}{2.3pt}}\par\vspace{1.15cm}
{\fontsize{24}{30}\selectfont\bfseries\color{aeNavy}\raggedright #1\par}\vspace{0.45cm}
{\Large\color{aeBlue}#2\par}\vspace{1.3cm}
{\large\bfseries #3\par}\vspace{0.35cm}
{\color{aeGray}Edition 1.0 \quad | \quad 29 October 2026\par}\vfill
{\color{aeGray}\small Original educational material aligned to the chapter and section topics in the five supplied volumes of Donald E. Knuth's \emph{The Art of Computer Programming}. This document is not an authorized reproduction or substitute for the original books.}\par\vspace{0.4cm}
{\color{aeTeal}\rule{\linewidth}{0.75pt}}
\end{titlepage}\hypersetup{pageanchor=true}\pagenumbering{roman}}
\newcommand{\SourceNote}[1]{\par\smallskip{\footnotesize\color{aeGray}\textbf{Source locator:} #1}\par\smallskip}
\newcommand{\Outcome}[1]{\par\textbf{Competency target.} #1\par}
\newcommand{\Checkpoint}[1]{\par\textbf{Checkpoint.} #1\par}
% Intentionally simple printable writing lines: call \AnswerLines{number}.
\newcount\alcount
\newcommand{\AnswerLines}[1]{\par\alcount=0\loop\ifnum\alcount<#1\par\noindent\rule{\linewidth}{0.3pt}\vspace{0.33cm}\advance\alcount by1\repeat\par}
\newcommand{\FormField}[2][0.62\linewidth]{\noindent\textbf{#2}\hspace{0.65em}\rule{#1}{0.4pt}\par\vspace{0.18cm}}
\newcommand{\ReviewTable}{\begin{tabularx}{\linewidth}{@{}p{0.23\linewidth}X p{0.2\linewidth}@{}}\toprule Criterion & Recorded evidence & Rating (0--4)\\\midrule Correctness & & \\[1em] Proof quality & & \\[1em] Complexity & & \\[1em] Testing & & \\[1em] Reproducibility & & \\[1em] Communication & & \\[1em]\bottomrule\end{tabularx}}

\begin{document}
\DocTitle{Volume 3 Study Guide}{Chapters 5--6: ordering, retrieval and storage}{Chapter- and section-aligned study guide}
\tableofcontents
\clearpage
\pagenumbering{arabic}
\section{Study protocol}

For each subsection, read the specified part of the original volume, keep a mathematical notebook, produce an original implementation or design sketch, and evaluate the acceptance gate. The investigations below are new exercises, not copies of the book's problems or solutions. Annotate any starred or mathematically advanced source sections as optional on the first pass; revisit them after the prerequisite mathematics is in place.
\begin{evidencebox}Recommended artifacts: study notes, problem specification, derivation, runnable code or model, tests, performance evidence and reflective comparison to original assumptions.\end{evidencebox}

\section{\S 5.1: Permutation statistics}
\SourceNote{Knuth, Volume 3, \S 5.1.}
\Outcome{Use inversion counts, runs and multiset permutations to characterize workload shape.}
\subsection{Independent investigation}
Generate arrays with controlled inversion counts and compare adaptive sorting behavior.
\subsection{Analysis and verification gate}
State how duplicate keys change the combinatorial model.
\subsection{Study questions}
\begin{enumerate}
\item Which representations, input assumptions and cost metrics control this technique's applicability?
\item Which invariant or combinatorial identity establishes correctness, and which input is its strongest counterexample candidate?
\item What result would cause the implementation to fail review despite passing a small collection of examples?
\end{enumerate}
\Checkpoint{Preserve a short explanation, the relevant test evidence, and one unresolved question before proceeding.}

\section{\S 5.2.1--5.2.3: Insertion, exchange and selection sorting}
\SourceNote{Knuth, Volume 3, \S 5.2.1--5.2.3.}
\Outcome{Analyze stability, comparison count and data movement for elementary sorting families.}
\subsection{Independent investigation}
Implement insertion sort and selection sort with key-plus-original-index records.
\subsection{Analysis and verification gate}
Show which implementations preserve equal-key order and why.
\subsection{Study questions}
\begin{enumerate}
\item Which representations, input assumptions and cost metrics control this technique's applicability?
\item Which invariant or combinatorial identity establishes correctness, and which input is its strongest counterexample candidate?
\item What result would cause the implementation to fail review despite passing a small collection of examples?
\end{enumerate}
\Checkpoint{Preserve a short explanation, the relevant test evidence, and one unresolved question before proceeding.}

\section{\S 5.2.4--5.2.5: Merging and distribution}
\SourceNote{Knuth, Volume 3, \S 5.2.4--5.2.5.}
\Outcome{Compare comparison-based merging with domain-restricted distribution sorting.}
\subsection{Independent investigation}
Implement stable merge sort and compare against a standard-library reference.
\subsection{Analysis and verification gate}
Test empty, singleton, duplicate, adversarial and nearly sorted inputs.
\subsection{Study questions}
\begin{enumerate}
\item Which representations, input assumptions and cost metrics control this technique's applicability?
\item Which invariant or combinatorial identity establishes correctness, and which input is its strongest counterexample candidate?
\item What result would cause the implementation to fail review despite passing a small collection of examples?
\end{enumerate}
\Checkpoint{Preserve a short explanation, the relevant test evidence, and one unresolved question before proceeding.}

\section{\S 5.3: Optimum sorting and lower bounds}
\SourceNote{Knuth, Volume 3, \S 5.3.}
\Outcome{Use decision-tree arguments to explain the $\Omega(n\log n)$ comparison lower bound.}
\subsection{Independent investigation}
Draw a decision tree for sorting three distinct elements and count possible leaves.
\subsection{Analysis and verification gate}
Explain why counting sort does not violate a comparison-model lower bound.
\subsection{Study questions}
\begin{enumerate}
\item Which representations, input assumptions and cost metrics control this technique's applicability?
\item Which invariant or combinatorial identity establishes correctness, and which input is its strongest counterexample candidate?
\item What result would cause the implementation to fail review despite passing a small collection of examples?
\end{enumerate}
\Checkpoint{Preserve a short explanation, the relevant test evidence, and one unresolved question before proceeding.}

\section{\S 5.4: External sorting}
\SourceNote{Knuth, Volume 3, \S 5.4.}
\Outcome{Model bounded memory, runs, multiway merge and I/O passes.}
\subsection{Independent investigation}
Design a k-way merge for a dataset larger than assigned RAM budget; create run size and fan-in calculations.
\subsection{Analysis and verification gate}
Compute bytes moved and include failure handling and temporary disk capacity.
\subsection{Study questions}
\begin{enumerate}
\item Which representations, input assumptions and cost metrics control this technique's applicability?
\item Which invariant or combinatorial identity establishes correctness, and which input is its strongest counterexample candidate?
\item What result would cause the implementation to fail review despite passing a small collection of examples?
\end{enumerate}
\Checkpoint{Preserve a short explanation, the relevant test evidence, and one unresolved question before proceeding.}

\section{\S 6.1--6.2.1: Sequential and ordered-table searching}
\SourceNote{Knuth, Volume 3, \S 6.1--6.2.1.}
\Outcome{Distinguish scans, lower bounds, binary search and index preconditions.}
\subsection{Independent investigation}
Implement lower-bound search and cross-check every insertion position on short sequences.
\subsection{Analysis and verification gate}
Prove maintenance of candidate interval and finite termination.
\subsection{Study questions}
\begin{enumerate}
\item Which representations, input assumptions and cost metrics control this technique's applicability?
\item Which invariant or combinatorial identity establishes correctness, and which input is its strongest counterexample candidate?
\item What result would cause the implementation to fail review despite passing a small collection of examples?
\end{enumerate}
\Checkpoint{Preserve a short explanation, the relevant test evidence, and one unresolved question before proceeding.}

\section{\S 6.2.2--6.2.4: Binary, balanced and multiway trees}
\SourceNote{Knuth, Volume 3, \S 6.2.2--6.2.4.}
\Outcome{Compare expected shape and guaranteed height and understand page-oriented search trees.}
\subsection{Independent investigation}
Sketch how update workload and storage medium change search-index choice.
\subsection{Analysis and verification gate}
Account for rotations/maintenance and for cache or disk block constraints.
\subsection{Study questions}
\begin{enumerate}
\item Which representations, input assumptions and cost metrics control this technique's applicability?
\item Which invariant or combinatorial identity establishes correctness, and which input is its strongest counterexample candidate?
\item What result would cause the implementation to fail review despite passing a small collection of examples?
\end{enumerate}
\Checkpoint{Preserve a short explanation, the relevant test evidence, and one unresolved question before proceeding.}

\section{\S 6.3--6.4: Digital searching and hashing}
\SourceNote{Knuth, Volume 3, \S 6.3--6.4.}
\Outcome{Evaluate tries, collision policies, load factor and hash-function assumptions.}
\subsection{Independent investigation}
Implement a simple collision-count simulator at several table load factors.
\subsection{Analysis and verification gate}
Explain how adversarial keys undermine expected-time assumptions.
\subsection{Study questions}
\begin{enumerate}
\item Which representations, input assumptions and cost metrics control this technique's applicability?
\item Which invariant or combinatorial identity establishes correctness, and which input is its strongest counterexample candidate?
\item What result would cause the implementation to fail review despite passing a small collection of examples?
\end{enumerate}
\Checkpoint{Preserve a short explanation, the relevant test evidence, and one unresolved question before proceeding.}

\section{\S 6.5: Secondary-key retrieval}
\SourceNote{Knuth, Volume 3, \S 6.5.}
\Outcome{Understand materialized indexes and multi-attribute query trade-offs.}
\subsection{Independent investigation}
Design two indexes for a synthetic record store and benchmark selective and broad predicates.
\subsection{Analysis and verification gate}
Report maintenance, storage and query overhead for each choice.
\subsection{Study questions}
\begin{enumerate}
\item Which representations, input assumptions and cost metrics control this technique's applicability?
\item Which invariant or combinatorial identity establishes correctness, and which input is its strongest counterexample candidate?
\item What result would cause the implementation to fail review despite passing a small collection of examples?
\end{enumerate}
\Checkpoint{Preserve a short explanation, the relevant test evidence, and one unresolved question before proceeding.}

\section{Volume synthesis and mastery review}
\begin{enumerate}
\item Draw a dependency graph connecting at least four major ideas from this volume.
\item Select two alternative implementations for one workload and compare their preconditions, proven cost, memory use, measured performance and failure cases.
\item Present a five-minute technical defense of the most surprising experimental result. A claim that differs from the mathematical model requires a measured or documented explanation.
\end{enumerate}
\subsection{Completion checklist}
\begin{itemize}\item All module notes include a source locator and an independently stated learning outcome.
\item Every implementation has reproducible tests and edge cases.
\item Every complexity claim identifies its cost model and underlying assumptions.
\item A peer can independently reproduce one chosen experiment using its recorded configuration.
\end{itemize}
\end{document}
